\documentclass[border=5pt]{standalone}
\usepackage[T1]{fontenc}
\usepackage{tikz}
\usepackage{luacode}

\definecolor{horseshoe}{HTML}{65780B}
\definecolor{atlantic}{HTML}{466A9F}
\definecolor{garnet}{HTML}{73000A}
\definecolor{rose}{HTML}{CC2E40}
\definecolor{warmgrey}{HTML}{676156}
\definecolor{honeycomb}{HTML}{A49137}
\definecolor{congaree}{HTML}{1F414D}
\definecolor{bk70}{HTML}{5C5C5C}
\definecolor{bk30}{HTML}{C7C7C7}

% Camera icon macro
\newcommand{\drawcameraicon}[1][1]{%
  \begin{scope}[scale=#1]
    \fill[bk70] (-0.18,-0.55) rectangle (0.18,0.55);
    \draw[black, thick] (-0.18,-0.55) rectangle (0.18,0.55);
    \fill[bk30] (0.18,-0.25) rectangle (0.38,0.25);
    \draw[black, thick] (0.18,-0.25) rectangle (0.38,0.25);
    \fill[bk30] (0.38,-0.33) rectangle (0.5,0.33);
    \draw[black, thick] (0.38,-0.33) rectangle (0.5,0.33);
    \fill[rose] (0.0,-0.38) circle (0.07);
    \draw[black, thin] (0.0,-0.38) circle (0.07);
  \end{scope}%
}

% ── Light state computed from \PARAM in Lua ──
% \PARAM ranges 0-94 (95 frames at 30fps)
% 0-4: all off, 5-34: singles, 35-64: pairs, 65-94: triplets

% ── Object ──
\def\objx{2.0}
\def\objy{0}
\def\objr{0.8}

% ── Camera ──
\def\camx{-2.2}
\def\camy{0}

\begin{document}
\begin{tikzpicture}
  \useasboundingbox (-3,-3) rectangle (7,3);
  \draw[horseshoe, thick] (-3,-3) rectangle (7,3);

\begin{luacode}
local objx, objy, objr = \objx, \objy, \objr
local camx, camy = \camx, \camy
local cam_radius = 0.15

-- 10 lights equally spaced in a ring around the object
-- Ring radius, skip the angle closest to camera (left side)
local nlights = 30
local ring_r = 2.2
local light_positions = {}
local palette = {"garnet", "rose", "atlantic", "congaree", "horseshoe", "honeycomb"}
local light_colors = {}
for i = 1, nlights do
  table.insert(light_colors, palette[((i-1) % #palette) + 1])
end

for i = 0, nlights - 1 do
  -- Start from right (0 deg), go counterclockwise
  -- Skip camera direction (~180 deg) by offsetting start
  local ang = (i / nlights) * 2 * math.pi + 0.3  -- offset to avoid camera
  local lx = objx + ring_r * math.cos(ang)
  local ly = objy + ring_r * math.sin(ang)
  table.insert(light_positions, {x=lx, y=ly, ang=ang})
end

-- Compute light states from frame number
local frame = math.floor(\PARAM + 0.5)
local states = {}
for i = 1, nlights do states[i] = 0 end

if frame < 5 then
  -- All off
elseif frame < 35 then
  -- Single light: light (frame-5)+1
  local idx = ((frame - 5) % nlights) + 1
  states[idx] = 1
elseif frame < 65 then
  -- Adjacent pairs
  local idx = ((frame - 35) % nlights) + 1
  states[idx] = 1
  states[(idx % nlights) + 1] = 1
else
  -- Adjacent triplets
  local idx = ((frame - 65) % nlights)
  states[(idx % nlights) + 1] = 1
  states[((idx + 1) % nlights) + 1] = 1
  states[((idx + 2) % nlights) + 1] = 1
end

local nrays = 100

-- Draw rays for active lights
for li, light in ipairs(light_positions) do
  local is_on = states[li] or 0
  if is_on == 1 then
    local dx = objx - light.x
    local dy = objy - light.y
    local dist = math.sqrt(dx*dx + dy*dy)
    local center_ang = math.atan2(dy, dx)
    local half_ang = math.asin(math.min(objr / dist, 1)) * 1.5

    local clr = light_colors[li] or "rose"

    for i = 0, nrays - 1 do
      local frac = i / (nrays - 1)
      local ang = center_ang - half_ang + frac * 2 * half_ang
      local rdx = math.cos(ang)
      local rdy = math.sin(ang)

      local ox = light.x - objx
      local oy = light.y - objy
      local a = rdx*rdx + rdy*rdy
      local b = 2*(ox*rdx + oy*rdy)
      local c = ox*ox + oy*oy - objr*objr
      local disc = b*b - 4*a*c

      if disc >= 0 then
        local t1 = (-b - math.sqrt(disc)) / (2*a)
        local t2 = (-b + math.sqrt(disc)) / (2*a)
        local t = t1
        if t < 0.01 then t = t2 end

        if t > 0.01 then
          local hx = light.x + t * rdx
          local hy = light.y + t * rdy

          local nx = (hx - objx) / objr
          local ny = (hy - objy) / objr

          local dot = rdx*nx + rdy*ny
          local refx = rdx - 2*dot*nx
          local refy = rdy - 2*dot*ny

          local ref_len = math.sqrt(refx*refx + refy*refy)
          refx = refx / ref_len
          refy = refy / ref_len

          local to_cam_x = camx - hx
          local to_cam_y = camy - hy
          local proj = to_cam_x * refx + to_cam_y * refy
          local closest_x = hx + proj * refx
          local closest_y = hy + proj * refy
          local miss_dist = math.sqrt((closest_x - camx)^2 + (closest_y - camy)^2)
          local hits_cam = (proj > 0) and (miss_dist < cam_radius * 3)

          -- Incoming ray
          tex.sprint(string.format(
            "\\draw[%s, thin, opacity=0.25] (%.4f,%.4f) -- (%.4f,%.4f);",
            clr, light.x, light.y, hx, hy
          ))

          if hits_cam then
            tex.sprint(string.format(
              "\\draw[atlantic, thin, opacity=0.6] (%.4f,%.4f) -- (%.4f,%.4f);",
              hx, hy, camx, camy
            ))
          else
            local ref_end_x = hx + refx * 1.5
            local ref_end_y = hy + refy * 1.5
            tex.sprint(string.format(
              "\\draw[%s, thin, opacity=0.06] (%.4f,%.4f) -- (%.4f,%.4f);",
              clr, hx, hy, ref_end_x, ref_end_y
            ))
          end
        end
      end
    end
  end
end

-- Check if any light is on
local any_on = false
for _, s in ipairs(states) do
  if s == 1 then any_on = true; break end
end

-- Draw light source markers
for li, light in ipairs(light_positions) do
  local is_on = states[li] or 0
  local clr = light_colors[li] or "rose"
  if is_on == 1 then
    -- Active light: always show
    tex.sprint(string.format(
      "\\fill[%s] (%.4f,%.4f) circle (0.1);",
      clr, light.x, light.y
    ))
    tex.sprint(string.format(
      "\\node[font=\\tiny\\sffamily, color=%s] at (%.4f,%.4f) {%d};",
      clr, light.x, light.y + 0.22, li
    ))
  elseif not any_on then
    -- All off: show all lights as hollow
    tex.sprint(string.format(
      "\\draw[%s, thin] (%.4f,%.4f) circle (0.1);",
      clr, light.x, light.y
    ))
    tex.sprint(string.format(
      "\\node[font=\\tiny\\sffamily, color=%s!30] at (%.4f,%.4f) {%d};",
      clr, light.x, light.y + 0.22, li
    ))
  end
  -- If any_on and this light is off: don't draw it
end
\end{luacode}

  % ── Object (drawn on top) ──
  \fill[warmgrey!25] (\objx,\objy) circle (\objr);
  \draw[warmgrey, thick] (\objx,\objy) circle (\objr);

  % ── Camera ──
  \pgfmathsetmacro{\camaimang}{atan2(\objy-\camy, \objx-\camx)}
  \begin{scope}[shift={(\camx,\camy)}, rotate=\camaimang]
    \drawcameraicon[0.55]
  \end{scope}

\end{tikzpicture}
\end{document}
